% =====================================================================
% paper06_assessment_separation_v75
% Parent: Paper06 v74 (preserved unchanged); it inherits the v73 Figure 2 correction.
% v75 (2026-10-27): source-checked, Gemini-based prose revision of §§1.1–1.2.
% The model, parameters, numerical results, theorem statements, and proofs are unchanged.
% Build: tectonic paper06_assessment_separation_v75.tex (elsarticle, authoryear)
% =====================================================================
\documentclass[preprint,11pt,authoryear]{elsarticle}

\usepackage{amsmath,amssymb,amsthm}
\usepackage{booktabs}
\usepackage{array}
\usepackage{tabularx}
\usepackage{graphicx}
\graphicspath{{./}{figs/}{paper06_assessment_separation_v75_src/figs/}}
\usepackage{enumitem}
\usepackage[colorlinks=true,urlcolor=blue,linkcolor=blue,citecolor=blue]{hyperref}
\emergencystretch=3em
\setlength{\tolerance}{2000}

\newtheorem{theorem}{Theorem}
\newtheorem{proposition}{Proposition}
\newtheorem{lemma}{Lemma}
\newtheorem{corollary}{Corollary}
\theoremstyle{definition}
\newtheorem{definition}{Definition}
\theoremstyle{remark}
\newtheorem{remark}{Remark}

\newcommand{\Etyp}{E_{\mathrm{typ}}}
\newcommand{\Eend}{E_{\mathrm{end}}}
\newcommand{\Ephys}{E_{\mathrm{tube,phys}}}
\newcommand{\Vtyp}{\mathcal{V}_{\mathrm{typ}}}
\newcommand{\Vweak}{\mathcal{V}_{\mathrm{weak}}}
\newcommand{\Vphys}{\mathcal{V}_{\mathrm{phys}}}
\newcommand{\Vend}{\mathcal{V}_{\mathrm{end}}}
\newcommand{\Wp}{W_+}
\newcommand{\FP}{\mathrm{FP}_{\mathrm{agg}}}
\newcommand{\Tube}{\mathrm{Tube}}
\newcommand{\Succ}{\mathrm{Succ}}
\newcommand{\End}{\mathrm{End}}
\newcommand{\conv}{\operatorname{conv}}
\newcommand{\Rp}{\mathbb{R}^n_+}
\newcommand{\Q}{\mathcal{Q}}
\newcommand{\act}[1]{\textsc{#1}}

\journal{Journal of Multi-Criteria Decision Analysis}

\begin{document}

\begin{frontmatter}

\title{Transition Certification with Composite Indices: The Acceptance Gap Between Compensatory and Noncompensatory Rules}

\author[1]{Amin Abaee\corref{cor1}}
\ead{amin_abaee@ut.ac.ir}
\cortext[cor1]{Corresponding author. ORCID: 0000-0002-0019-1842.}
\affiliation[1]{organization={Independent Researcher}, city={Tehran}, country={Iran}}

\begin{abstract}
Composite indices certify transitions under weight-adaptive or common-plan rules. We compare them as robust finite-horizon problems with path-wise constraints and disturbances. The full-cone aggregate is lossless at a fixed trajectory, yet the plan/weight quantifiers need not commute. For finite menus of required-margin vectors, the gap is $\conv(\mathcal D)+\Rp$ outside the union of single-plan acceptance orthants; convexifying the menu closes it.

An exact rational witness has a nonempty gap interior: at its canonical point both floors are $6/5$ and the equal-weight aggregate $12/5$ clears $2$, though no primitive plan meets both. A closed-form reserve threshold separates rescuable from impossible states; exact weight ratios identify the licensed plan on the gap interior. Blending closes the action-indexed gap, not the shared-shock gap; full-depth alternation does not implement a blend. On the witness-defined diagonal gap slice, every positive-elasticity CES index falsely certifies an exact interval; the Leontief limit is typed-safe. The shared-shock gap survives the exact $1/2$-cut; the distinct replace reading with identical plan- and weight-independent dips collapses the distinction (SI S4). Published DFO biomass/LRP figures are a limited scale/context anchor. Per-floor path minima, reserve and protocol reconstruct the typed verdict; the aggregate alone does not.
\end{abstract}

\begin{keyword}
multi-criteria decision analysis \sep compensatory and noncompensatory aggregation \sep composite indicators \sep scalarization \sep robustness \sep sustainability assessment \sep viability
\end{keyword}

\end{frontmatter}

\section{Introduction}\label{sec:intro}

\subsection{Compensability as an acceptance question}

Whenever decision-makers assess a sustainability transition, they face a familiar question: can strength in one area make up for failure in another? Multi-attribute value theory builds trade-offs into its value structure \citep{KeeneyRaiffa1976}. Outranking methods, by contrast, limit compensation through concordance and veto conditions \citep{Roy1996,Vincke1992}; work on sustainable development has also emphasized the incommensurability of some criteria \citep{Munda2005,MartinezAlier1998}.

The same question appears in environmental assessment as the distinction between weak and strong sustainability. \emph{Weak sustainability} represents substitution among forms of capital through an aggregate that must not decline \citep{WorldBank2011,Neumayer2013,Boos2015}. \emph{Strong sustainability} treats selected stocks as separately binding \citep{Daly1990,Ekins2003}. Research on the indicator theory of these positions is well developed, as is the analysis of compensability in the multi-criteria methods used for sustainability assessment \citep{Hanley1999,Martinet2011,CairnsMartinet2014,DoyenGajardo2020,Usubiaga2025,Cinelli2014,Schar2025}.

Much of this work classifies aggregation at a fixed profile of criterion values—a snapshot. This paper asks a different question, about \emph{transitions}: when a plan unfolds over time under declared disturbances, and several criteria must hold along the entire path, when can a compensatory aggregate certify a transition that a noncompensatory assessment rejects? The question is not how weights should be chosen; methods for eliciting them are well developed \citep{Nardo2008,Becker2017}. Nor is it the familiar static limitation of weighted sums in recovering nonconvex parts of a Pareto front \citep{DasDennis1997}. The focus is \emph{acceptance semantics}: the rule that classifies a composite reading of a transition as accepted or rejected—the sorting problem of decision aiding \citep{Roy1996}.

\subsection{Two acceptance semantics}

To see where the two assessments can diverge, imagine the transition problem at state $z$. It tracks typed-margin coordinates $s_i$, normalised so that zero is each binding floor, and has a finite menu $\mathcal A$ of plans and a declared disturbance set $D$. Both assessments impose the same physical safe and destination sets $S^{\mathrm{phys}}$ and $G^{\mathrm{phys}}$. Exact-tube semantics checks the entire path, not just its endpoint: for every $d\in D$, a plan's full tube must remain in $S^{\mathrm{phys}}$, and its successor set must lie in $G^{\mathrm{phys}}$.

\begin{samepage}
Now consider the floor test. Write $s(p)=(s_1(p),\dots,s_n(p))$ for the typed-margin vector at state $p$. For each weight $w$ in the full nonnegative cone $\Wp=\Rp\setminus\{0\}$, let $E_w(z)$ be the menu plans that meet those shared requirements and also keep $w\cdot s(p)\ge0$ at every tube state and every successor state, for every $d\in D$. Let $\Etyp(z)$ be the menu plans that meet the same physical, destination, and disturbance requirements but keep each typed floor $s_i(p)\ge0$ at those same states. Thus the two assessments share physical feasibility, destination attainment, and disturbance coverage; only the test applied to the typed margins differs.
\end{samepage}

The key choice is simple: must one plan work for every weighting, or may each weighting have its own plan? Formally, the two acceptance rules differ only in the order of the quantifiers:
\begin{itemize}[leftmargin=2em]
\item \emph{Common-plan acceptance}: $\exists a\in\mathcal A\;\forall w\in\Wp:\ a\in E_w(z)$. One plan must work for every weighting.
\item \emph{Weight-adaptive acceptance}: $\forall w\in\Wp\;\exists a_w\in\mathcal A:\ a_w\in E_w(z)$. For each weighting, some plan works; it may be a different one.
\end{itemize}
The first rule implies the second, but the converse need not hold. Because $\Wp$ contains every coordinate weight, the full-cone identity is $\Etyp(z)=\bigcap_{w\in\Wp}E_w(z)$. Among plans meeting the shared requirements, a single plan passes every weight exactly when it satisfies each typed floor; equivalently, common-plan acceptance is $\exists a\in\Etyp(z)$. The gap, when it occurs, comes from the weight-adaptive rule's ability to use different plans for different weights. Table~\ref{tab:protocols} places these two rules among the four protocols considered in the paper.

\begin{table}[!htbp]
\centering\small
\caption{Four acceptance protocols for a composite index on a transition datum. Protocols 3 (common-plan robustness) and 4 (typed acceptance) have the same accepted-plan and accepted-state sets on the full cone (Proposition~\ref{prop:hierarchy}(ii)); they are shown separately to distinguish the common-plan and typed formulations. The separation is Protocol 2 against either equivalent form.}
\label{tab:protocols}
\begin{tabularx}{\textwidth}{>{\raggedright\arraybackslash}p{0.25\textwidth} >{\raggedright\arraybackslash}p{0.27\textwidth} X}
\toprule
Protocol & Quantifier form & Reading \\
\midrule
1. Fixed-weight assessment & $a\in E_w(z)$ at a declared $w$ & one index, one selected plan: ordinary fixed-weight certification \\
2. Weight-adaptive (per-weight) acceptance & $\forall w\ \exists a_w\in E_w(z)$ & for each weighting $w$, some plan $a_w$ is licensed; the plan may vary with $w$: the compensatory reading, $\Vweak=\bigcap_w\mathcal V_w$ \\
3. Common-plan robustness & $\exists a\ \forall w:\ a\in E_w(z)$ & one plan serves every weighting; equivalent to Protocol 4 on the full cone \\
4. Coordinate-wise (typed) acceptance & $\exists a:\ a\in\Etyp(z)$ & separately binding floors enforced path-wise; equivalent to Protocol 3 on the full cone \\
\bottomrule
\end{tabularx}
\end{table}

An analogous split appears in adjustable robust optimisation: a here-and-now decision is fixed before the relevant parameter is known, while a wait-and-see decision may depend on it. Here the assessment weight plays the role of that parameter \citep{BenTal2004}. This is an analogy for the order of the acceptance rule, not a claim about how real auditors or policymakers behave. The result concerns a certification protocol applied to a composite index, not the suitability of composite indices as descriptive summaries.

The witness makes the split concrete. \act{fast} risks the first floor, \act{slow} the second. \act{staged} protects both only when the reserve covers its bridge; \act{no-switch} misses the destination. At the action-indexed canonical point $(x,s_1,s_2)=(\tfrac12,\tfrac65,\tfrac65)$, $s_1+s_2=\tfrac{12}5\ge2$, so every $w\in\Wp$ licenses \act{fast} or \act{slow}, yet none of the four plans is typed-admissible. This verdict is specific to the stated menu and disturbance class.


\pagebreak
\subsection{Contributions}

The paper makes five contributions, stated in the order of the text.
\begin{enumerate}[leftmargin=2em,label=(\roman*)]
\item \emph{Localisation and hierarchy} (Section~\ref{sec:general}). On the full cone the typed operator is exactly the common-plan intersection, $\Etyp(z)=\bigcap_{w\in\Wp}E_w(z)$, and at any fixed trajectory the aggregate is lossless (Lemma~\ref{lem:cone}); the chain typed $\subseteq$ scalarized $\subseteq$ physical $\subseteq$ endpoint holds for every weight, and restricting the weight family can only enlarge the compensatory accepted set (Propositions~\ref{prop:hierarchy}--\ref{prop:monotone}). The separation is therefore dynamic -- a property of the quantifier over plans -- and the full cone is its strictest case.
\item \emph{A geometric characterisation for any finite menu} (Theorem~\ref{thm:menu}). For plans whose coordinate-wise admissibility is governed by a required-margin vector $d_a$, weight-adaptive acceptance is $\conv(\mathcal D)+\Rp$ and common-plan acceptance is $\bigcup_a(d_a+\Rp)$; the acceptance gap is the part of the up-closed convex hull $\conv(\mathcal D)+\Rp$ not covered by any single plan, and convexifying the menu closes it. A synchronisation lemma (Lemma~\ref{lem:sync}) fixes the scope: the margin-vector geometry describes the protocol that scalarizes within-path maxima exactly, and is conservative for a protocol that scalarizes first and maximises afterwards.
\item \emph{An explicit rational witness} (Theorem~\ref{thm:witness}). On a four-plan datum the gap is a region with nonempty interior, equal to the \emph{impossibility region} $I$; the aggregate-versus-floor discrepancy region splits at the reserve level $x=1$ into $I$ and a \emph{rescue set} $R$ that a reserve-financed staged plan already serves; exact thresholds $\rho_1=(2-s_1)/s_2$ and $\rho_2=s_1/(2-s_2)$ locate the licensing regions of the two principal plans, with coordinate-axis weights classified separately; and the minimal reserve increment that converts impossibility into rescue is $\kappa^*(z)=(1-x)\,\mathbf 1[z\notin\Vtyp]$ (Proposition~\ref{prop:rescue}). The separation persists under hold-prefix extension of the horizon and can be erased by a later stage with coincident terminal sets (Proposition~\ref{prop:persist}). The quantifier gap is sufficient but not necessary for false acceptance: at the canonical point $(x,s_1,s_2)=(1/2,6/5,6/5)$ the single fixed weight $(1,1)$ already accepts a primitive plan, because $12/5\ge2$, while no primitive plan meets both typed floors (Remark~\ref{rem:fixedweight}).
\item \emph{What closes the gap and what does not} (Theorem~\ref{thm:blend}, Proposition~\ref{prop:timeshare}). Fractional blending of the primitive plans -- menu convexification at the level of control flows -- makes the typed accepted set of the action-indexed margin class coincide exactly with the compensatory region, so that gap is a property of the finite deterministic menu, not of either doctrine. The statement is class-specific: under the shared-shock class the blend window is empty at the canonical point ($12/5<5/2$), so blending does not close that gap (Proposition~\ref{prop:commonshock}). Full-depth alternation between the primitive plans does not implement a blend and closes the gap only where both dips are separately survivable; arbitrarily fine duty cycling is a different control class, which the alternation result does not rule out.
\item \emph{Robustness of the separation to the aggregator and to the disturbance class} (Theorem~\ref{thm:ladder}, Proposition~\ref{prop:commonshock}). Replacing the linear index by a constant-elasticity-of-substitution (CES) aggregator with elasticity $\sigma$ changes the gap but does not remove it: the accepted rungs are nested, on the witness-defined diagonal gap slice, each positive-elasticity member falsely certifies an exact interval with critical floors $1<5/4<\sqrt2<\varphi<\sqrt3<\dots\to2$; therefore no positive elasticity is uniformly safe over the full gap, while the Leontief member ($\sigma=0$), which is the typed operator, is uniformly safe. Under a shared environmental shock striking every plan the gap survives, reduced by an exact $1/2$-cut, with the licensing thresholds unshifted on the surviving region; the reserve route then fails on an exact band of states.
\end{enumerate}
Every continuum statement is proved (\ref{app:general}--\ref{app:witness}; Supplementary Information, SI~S2); the exact-arithmetic check counts and the scope of the artifacts available for independent reproduction are recorded in SI~S7 and Data Availability. The core theorems are additionally formalised in Lean~4 (Section~\ref{sec:verification}). A resource-transition benchmark (Section~\ref{sec:benchmark}) realises the datum as a fishery transition with units anchored to a regulated biomass limit.

\paragraph{What is not claimed} No ranking of weak and strong sustainability as doctrines; no proof that any particular set of boundaries ought to be treated noncompensatorily; no empirical fit, forecast, validation, or performance result for a resource system (published DFO values are a limited scale/context anchor only); no infinite-horizon, stochastic or partially observed extension. The impossibility region is a negative certificate relative to the specified plan menu, and the separation is exhibited for a declared disturbance class.

\subsection{Position in the literature}

The backward constructions used here are typed instances of robust-predecessor, capture-basin and reach--avoid constructions \citep{Aubin1991,Aubin2011,SaintPierre1994,Lygeros1999,AubinCatte2002}; the hierarchy of Proposition~\ref{prop:hierarchy}(i) is constraint-set monotonicity of viability kernels \citep{Aubin1991,Frankowska1989}; viability methods are established in natural-resource management \citep{DeLaraDoyen2008,Doyen2012}. The blend-collapse pair is the pure-versus-mixed-strategy structure of the assessor--planner game \citep{vonNeumann1928,Sion1958}, whose reach--avoid form is the discriminating kernel of \citet{Cardaliaguet1996}. In composite-indicator practice the critique of weights as value judgements presented as measurement \citep{Hickel2020,Fischer2022} and the construction of indices without weights from Pareto fronts \citep{Gao2023} both concern the index; the result here is orthogonal -- even with weights elicited faithfully, the quantifier order produces the gap. We do not assert priority over the elementary fact $\exists a\,\forall w\ne\forall w\,\exists a_w$, nor over the scalarization gap of multi-objective optimisation of which Theorem~\ref{thm:menu}(ii) is an instance; Section~\ref{sec:related} places the results against adjustable robust optimisation, the fixed-order and path-wise worst cases of multi-stage decision theory, randomisation-receptive assessment and the weak/strong sustainability debate, and states exactly which four claims are made as new and which are not.
